\section{Formal Framework}

\subsection{Typed relational carriers}

We use the category $\FinRel$ of finite sets and binary relations. Its objects are
finite sets, and a morphism $R:X\nrightarrow Y$ is a subset of $X\times Y$.
For $R:X\nrightarrow Y$ and $S:Y\nrightarrow Z$,
\[
S\circ R=\{(x,z):\exists y\in Y,
(x,y)\in R,\ (y,z)\in S\}.
\]
The identity on $X$ is $\Delta_X=\{(x,x):x\in X\}$.

Fix a finite quiver $Q=(Q_0,Q_1,s,t)$ and a finite set $\mathcal E$ of
pairs of parallel paths in its free path category. The category
$\operatorname{Path}(Q)$ has the vertices of $Q$ as objects, finite composable
paths as morphisms, length-zero paths as identities, and concatenation as
composition. The schema is the presented
category
\[
\mathcal S=\operatorname{Path}(Q)/\!\equiv_{\mathcal E},
\]
where $\equiv_{\mathcal E}$ is the smallest path congruence containing
$\mathcal E$. The arrows of $Q$ are the declared generators, and every schema
morphism has a path representative subject to the declared equations. A typed
relational carrier is a functor $F:\mathcal S\to\FinRel$ such that every
$F(A)$ is nonempty and finite. A generator $e\in Q_1$ is interpreted as the
relation $F([e])$. Functoriality requires
$F(\id_A)=\Delta_{F(A)}$ and $F(g\circ f)=F(g)\circ F(f)$.
A label family is a collection of maps
$\lambda_A:F(A)\to L_A$ into fixed typewise label sets.

\begin{definition}[Integrated structural imagination state]
An integrated structural imagination state is a tuple
\[
\Sigma=(F,\lambda,K,d,f,\mu),
\]
where $F$ is a typed relational carrier, $\lambda$ is a label family,
$K$ is a finite abstract simplicial complex whose vertex set is the tagged union
$|F|=\coprod_A F(A)$, and $d$ is a metric on $|F|$. The function
$f:K\to\mathbb R$ is a monotone filtration and
$\mu:|F|\to(0,1]$ is a full-support probability mass function. The state also
obeys the compatibility condition
\[
(x,y)\in F([e])\Longrightarrow\{x,y\}\in K\qquad(e\in Q_1).
\]
The class of such states is denoted by $\StrState(Q,\mathcal E)$.
\end{definition}

The common tagged carrier is a scope decision. It permits relations, simplices,
metric distances, filtration values, and mass to refer to the same typed
entities. The compatibility condition is the minimal cross-layer constraint
used in this paper: a declared generator relation must be visible as a
one-simplex, but no converse or metric compatibility is assumed.

\begin{definition}[Integrated structural isomorphism]
Let $\Sigma=(F,\lambda,K,d,f,\mu)$ and
$\Sigma'=(G,\nu,L,\rho,g,\eta)$. An integrated structural isomorphism
$\varphi:\Sigma\to\Sigma'$ is a family of bijections
$\varphi_A:F(A)\to G(A)$ such that:
\begin{enumerate}
\item $\nu_A(\varphi_Ax)=\lambda_A(x)$ for every $A,x$;
\item the tagged union $\bar\varphi:|F|\to|G|$ satisfies
$\bar\varphi[K]=L$;
\item $\rho(\bar\varphi x,\bar\varphi y)=d(x,y)$ for all $x,y$;
\item $g(\bar\varphi\sigma)=f(\sigma)$ for every $\sigma\in K$;
\item $\eta(\bar\varphi x)=\mu(x)$ for every $x\in|F|$;
\item for every generator $e\in Q_1$,
\[
((\varphi_{s(e)}\times\varphi_{t(e)})[F([e])])=G([e]).
\]
\end{enumerate}
We write $\Sigma\cong_{\mathrm{str}}\Sigma'$ when such a family exists.
\end{definition}

\begin{definition}[Tracked transition]
A tracked transition $p:\Sigma\rightharpoonup\Sigma'$ is a family
$p=(p_A)_A$ in which each $p_A$ is an injective map from a subset of $F(A)$
into $G(A)$ and
\[
\nu_A(p_Ax)=\lambda_A(x)\qquad(x\in\dom p_A).
\]
Its domain and image tagged unions are written $D_p$ and $I_p$. The identity is
the total identity family. If
$p:\Sigma_0\rightharpoonup\Sigma_1$ and
$q:\Sigma_1\rightharpoonup\Sigma_2$, their composite is the family
$(q_A\circ p_A)|_{p_A^{-1}(\dom q_A)}$.
\end{definition}

\begin{definition}[Tagged extension of a partial transition]
For a tracked transition $p=(p_A)_A$, its tagged extension $\bar p$ is the
map on $D_p=\coprod_A\dom(p_A)$ defined by $\bar p(x)=p_A(x)$ when
$x\in F(A)$. Its image is $I_p=\coprod_A\im(p_A)$.
\end{definition}

\begin{proposition}[The category of tracked transitions]
Integrated states and tracked transitions form a category
$\StrTrack(Q,\mathcal E)$.
\end{proposition}

\begin{proof}
The identity family is injective, label-preserving, and has the full carrier as
its domain. A composite is injective on the restricted domain. If $x$ is in
its domain, label preservation follows successively from $p$ and $q$. The
domain of a triple composite is the set of elements on which all three partial
maps are defined, and both parenthesizations have the same value there. Hence
composition is associative. Restricting either composite by an identity leaves
the same partial family, proving the unit laws.
\end{proof}

\begin{definition}[Action-conditioned transition]
Let $\mathcal A$ be an action set. An action-conditioned transition is a pair
$(a,p)$ with $a\in\mathcal A$ and $p$ a tracked transition. The action label is
an index on the transition, not additional data in the morphism of
$\StrTrack(Q,\mathcal E)$; consequently composition is defined on the tracked
maps and does not concatenate incompatible single labels.
\end{definition}
